\documentclass[11pt]{article}
\usepackage[margin=1in]{geometry}
\usepackage{booktabs}
\usepackage[table]{xcolor}
\usepackage{amsmath}
\usepackage{enumitem}
\usepackage{titlesec}
\usepackage[hidelinks]{hyperref}
\setlist[itemize]{leftmargin=1.4em,itemsep=2pt,topsep=3pt}
\titleformat{\section}{\large\bfseries}{\thesection.}{0.5em}{}
\titleformat{\subsection}{\normalsize\bfseries}{\thesubsection}{0.5em}{}
\newcommand{\cls}[1]{\textbf{#1}}
\title{\vspace{-1.5em}Feature documentation --- \textbf{4-class} ground-condition model\\[2pt]
\large continuous\_paved / modular\_paved / loose / grass}
\author{GroundWalks (WalkSensePlace)}
\date{}
\begin{document}
\maketitle
\vspace{-2em}

Every feature is computed \textbf{per step} (one gait cycle, from one total-pressure peak to the
next, $\sim$0.5--1.2\,s at 62.5\,Hz). Sensors used: 12 plantar pressure pads, a 3-axis
accelerometer, and a 3-axis gyroscope (the magnetometer is dead --- all zeros --- so it is
unused). Steps that span a surface change or a \texttt{transition}/\texttt{wait} tag are
excluded, so each step lies on a single surface.

Two conventions appear throughout:
\begin{itemize}
\item \textbf{``norm'' / session-normalised} --- divided by a within-recording scale (e.g.\ that
  step's mean total pressure) so the value doesn't scale with body weight or how tightly the
  insole was worn.
\item \textbf{amplitude-invariant} --- the metric is a \emph{ratio} or \emph{shape} number,
  unaffected by overall size.
\end{itemize}

The model uses \textbf{48 features} (a cross-user pruned subset of 30 gives equivalent grouped
performance --- see \S8). The person-confound features the pipeline computes but \textbf{drops}
from the model are listed in \S7. The ``F'' values are the ANOVA class-separation scores
\emph{on the 4-class data} --- higher means the feature better separates the four classes.

\section{The 4-class scheme (what changed from 3-class)}
This model extends the merged 3-class model (grass / loose / paved) by \textbf{splitting paved}
along an \emph{evenness} axis:
\begin{itemize}
\item \cls{continuous\_paved} --- concrete, asphalt, brick, exposed aggregate: smooth, continuous
  rigid surfaces.
\item \cls{modular\_paved} --- cobblestone and uneven / loose-jointed brick: rigid like concrete
  but \emph{bumpy and segmented}. This is the \textbf{new class}.
\item \cls{loose} --- dirt, compact ground, gravel, sand.
\item \cls{grass}.
\end{itemize}
Crucially, continuous-vs-modular is \emph{not} a material distinction (both are hard, and the
sensor is weak at material). It rides on the \textbf{regularity / micro-vibration} signal --- the
step-to-step foot-tilt variability, jerk, and spectral texture features below --- which is the
axis the sensor genuinely responds to. Cobblestone is hard \emph{and} uneven, so it should sit
between smooth pavement and the softer classes on both the hardness and evenness scales; the
per-surface scale positions (\S9) confirm this.

\begin{center}
\rowcolors{2}{gray!8}{white}
\begin{tabular}{lcc}
\toprule
\textbf{surface} & \textbf{hardness (0--100)} & \textbf{evenness (0--100)}\\
\midrule
continuous\_paved & 64.4 & 55.7\\
modular\_paved    & 56.6 & 50.7\\
loose            & 53.2 & 54.4\\
grass            & 38.7 & 33.8\\
\bottomrule
\end{tabular}
\end{center}

\noindent\textbf{Data-status caveat.} As of this run, \cls{modular\_paved} has only $\sim$215
steps from \textbf{3 recordings}, versus 8055 steps / 37 recordings for continuous\_paved.
Under grouped cross-validation (unseen recordings) modular\_paved F1 is $\sim$18\%, and its errors
scatter mostly into continuous\_paved (both hard). The \emph{features} are the right ones --- the
top separators are the evenness/vibration features --- but the class is \textbf{data-starved, not
mis-modelled}. The reliable fix is more cobblestone recordings from more people; until then, read
the \emph{grouped} modular\_paved score (not the in-sample one) and expect it to be noisy.

\section{Foot-tilt variability (the strongest features)}
From a ZUPT-gated orientation estimate: within each step, the instant of minimum gyro magnitude
(foot flat and momentarily still) is found, and the foot's tilt there is read from the
accelerometer (which then points along gravity). Two orthogonal tilt axes, A and B, are tracked.
\begin{itemize}
\item \texttt{orient\_rest\_a\_roll3\_std} (F$\approx$1450) --- and \texttt{orient\_rest\_b\_roll3\_std} (F$\approx$668) --- the \textbf{step-to-step variability} of resting foot tilt (rolling 3-step std of flat-foot tilt within each recording+foot sequence). High values mean the foot lands at a \emph{different} orientation each step --- the signature of an uneven surface, and exactly the cue that separates modular from continuous paved. These are the single most class-differentiable features and generalise across people. The absolute resting tilt is computed but \textbf{not} used (per-person mount offset); only its variability is kept.
\item \texttt{orient\_range\_a} (F$\approx$74) --- and \texttt{orient\_range\_b} --- the range (max$-$min) of foot tilt \emph{during a single stance} on each axis: how much the foot rocks while planted. A within-step complement to the step-to-step variability above.
\end{itemize}

\section{Accelerometer texture / impact (vibration and hardness)}
Computed on the within-step accelerometer signal. Spectral features run an FFT on the detrended
signal; jerk/zcr run on vertical (accel-Y). These carry most of the continuous-vs-modular signal.
\begin{itemize}
\item \texttt{accel\_jerk\_ratio} (F$\approx$751) --- RMS of the first difference of vertical acceleration $\div$ RMS of the signal (amplitude-invariant). How ``jerky''/high-frequency the strike is --- cobblestone transmits sharper, more irregular vibration than smooth pavement.
\item \texttt{accel\_hf\_frac} (F$\approx$238) --- fraction of accelerometer spectral power above 5\,Hz. Higher on firm/rigid, textured surfaces.
\item \texttt{accel\_spec\_centroid} (F$\approx$252) --- spectral centroid (centre-of-mass frequency) of the accelerometer power spectrum. Shifts up with harder, higher-frequency impacts.
\item \texttt{accel\_spec\_entropy} (F$\approx$218) --- normalised Shannon entropy of the accelerometer spectrum: tonal/peaky (low) vs broadband (high). A texture descriptor.
\item \texttt{accel\_zcr} (F$\approx$240) --- zero-crossing rate of vertical acceleration: a simple roughness/oscillation measure.
\item \texttt{accel\_mag\_strike} (F$\approx$67) --- peak accelerometer magnitude in the impact region (first third of the step): strike intensity.
\end{itemize}

\section{Gyroscope motion (foot rotation \& stability)}
\begin{itemize}
\item \texttt{gyro\_x\_std} (F$\approx$596) --- standard deviation of X-axis angular rate during the step: side-to-side foot rotation/wobble through stance. A strong discriminator.
\item \texttt{gyro\_x\_strike} (F$\approx$276) --- and \texttt{gyro\_y\_strike} (F$\approx$109), \texttt{gyro\_z\_strike} (F$\approx$80) --- angular rate on each axis at the strike instant: the rotational ``kick'' at contact.
\item \texttt{gyro\_z\_std} (F$\approx$243) --- and \texttt{gyro\_y\_std} --- angular-rate variability on the other two axes.
\item \texttt{gyro\_hf\_frac} (F$\approx$192) --- fraction of gyroscope spectral power above 5\,Hz (rotational high-frequency content, analogous to \texttt{accel\_hf\_frac}).
\item \texttt{gyro\_mag\_strike} (F$\approx$195) --- gyro magnitude at the strike instant.
\end{itemize}

\section{Stance timing \& pressure-curve shape}
From the total-pressure curve over the step, mostly session-normalised so they describe
\emph{shape} not size.
\begin{itemize}
\item \texttt{stance\_duration} (F$\approx$349) --- time (samples) the foot is loaded. Compliant ground lengthens/softens contact.
\item \texttt{loading\_rate\_norm} (F$\approx$61) --- and \texttt{loading\_rate} --- how fast pressure rises to peak (session-normalised and raw). Hard surfaces load faster.
\item \texttt{rise\_slope}, \texttt{fall\_slope}, \texttt{slope\_asym} --- slopes of the pressure rise and fall, and their asymmetry.
\item \texttt{time\_to\_peak\_frac} --- fraction of the step elapsed when total pressure peaks.
\item \texttt{gu\_dip\_depth} (F$\approx$101) --- depth of the mid-stance dip in the amplitude-normalised total-pressure curve, $1-\min(\text{midstance})/\max$. Grass shows the deepest dip.
\item \texttt{gu\_load\_first\_frac} (F$\approx$105) --- fraction of total loading in the first half of the step.
\item \texttt{gu\_push\_peak} (F$\approx$96) --- height of the push-off (second) peak in the normalised curve.
\item \texttt{gu\_load\_slope} and \texttt{gu\_unload\_slope} (F$\approx$138) --- max rising / min falling slope of the normalised total curve (ground-up complements to the slope features).
\item \texttt{gu\_stance\_frac} (F$\approx$119) --- fraction of the step above half the mean total pressure: how much of the cycle is solidly loaded.
\end{itemize}

\section{Weight distribution (heel vs forefoot)}
\begin{itemize}
\item \texttt{heel\_fore\_ratio} (F$\approx$158) --- mean heel pressure $\div$ mean forefoot pressure: where load sits front-to-back.
\item \texttt{heel\_impulse\_norm} (F$\approx$176) --- heel pressure integrated over the step, session-normalised (time-under-load at the heel).
\item \texttt{heel\_skew} (F$\approx$62) --- and \texttt{heel\_kurt} (F$\approx$51) --- skewness and kurtosis of the heel-pressure time series (heel loading profile shape).
\item \texttt{heel\_ripple\_frac} (F$\approx$172) --- fraction of high-frequency ``ripple'' in the heel signal (mortar-joint chatter).
\item \texttt{midstance\_diff} --- heel-minus-forefoot balance at the mid-stance (flat-foot) instant.
\item \texttt{gu\_fore\_early\_frac} (F$\approx$90) --- share of forefoot loading in early stance.
\item \texttt{gu\_fore\_heel\_toff} --- timing offset between the forefoot and heel pressure peaks.
\end{itemize}

\section{Pressure-map distribution \& dynamics (spatial, across 12 pads)}
Treats the 12 sensors as a spatial pressure map.
\begin{itemize}
\item \texttt{gu\_spatial\_entropy\_stance} (F$\approx$210) --- and \texttt{gu\_spatial\_entropy\_ff} (F$\approx$38) --- normalised Shannon entropy of the pressure distribution across the 12 pads (over stance, and at flat-foot). High = load spread evenly; low = concentrated on a few pads. Compliant surfaces spread load more.
\item \texttt{gu\_participation\_ff} (F$\approx$43) --- effective fraction of pads bearing load at flat-foot (inverse participation ratio).
\item \texttt{gu\_active\_sensors} --- count of pads loaded above 10\% of the peak pad at flat-foot: contact breadth.
\item \texttt{gu\_cop\_wander} (F$\approx$81) --- mean frame-to-frame change of the normalised pressure map through stance: how much the load shifts around the sole (a dynamic unevenness proxy --- relevant to cobblestone).
\item \texttt{gu\_press\_jitter} (F$\approx$88) --- mean per-pad temporal jitter (std of the 2nd difference) of pressure, session-normalised: small fluctuations from surface irregularity.
\item \texttt{stance\_duration\_roll3\_std} (F$\approx$49) --- with \texttt{loading\_rate\_roll3\_std} and \texttt{impact\_mag\_roll3\_std} (F$\approx$42) --- step-to-step variability (rolling 3-step std) of stance duration, loading rate, and impact: gait-consistency measures that rise on irregular ground.
\end{itemize}

\section{Cross-user pruned feature set}
Cross-user permutation testing ranks features by how much they help a \emph{held-out person}.
Pruning the 18 features with $\le 0$ cross-user importance (to 30 features) gives essentially
identical grouped macro-F1 (33.2\% $\to$ 33.0\%) with less overfitting risk --- recommended as
data grows. The top cross-user features are \texttt{orient\_rest\_a\_roll3\_std}, \texttt{accel\_jerk\_ratio}, \texttt{orient\_rest\_b\_roll3\_std}, and \texttt{accel\_spec\_centroid}.

\section{Dropped (computed but excluded from the model)}
Removed because cross-user permutation testing showed they act as a \textbf{person fingerprint}
(who is walking / how the insole is seated, not the surface) and hurt generalisation:
\begin{itemize}
\item \textbf{Per-sensor pressure ratios} --- \texttt{pressure\_01\_mean\_ratio} $\ldots$ \texttt{pressure\_12\_mean\_ratio}, plus \texttt{heel\_mean\_ratio}, \texttt{fore\_mean\_ratio}, \texttt{heel\_pmax\_ratio}, \texttt{fore\_pmax\_ratio}.
\item \textbf{Raw magnitudes} --- \texttt{total\_p\_mean}, \texttt{total\_pmax}, \texttt{impact\_mag}, \texttt{impact\_mag\_norm}, \texttt{heel\_impulse} (scale with body weight / donning tightness).
\end{itemize}

\section*{How to read the model, in one line}
The features that carry the surface signal --- and survive to a new person --- are
\textbf{step-to-step foot-tilt variability}, \textbf{accelerometer texture/jerk}, \textbf{gyro
rotation variability}, \textbf{stance timing}, and \textbf{pressure-map spread}. Together these
describe \emph{compliance} (soft vs firm), which is why grass-vs-hard is reliable, \emph{and}
\emph{evenness/regularity}, which is the axis the new \textbf{continuous-vs-modular-paved} split
depends on. That split is promising but currently data-limited: its verdict rests on collecting
more cobblestone walks from more people.
\end{document}
